\documentclass[11pt, a4paper, openany, oneside]{book}

% --- Universal Preamble & Font Configuration ---
\usepackage[a4paper, top=2.5cm, bottom=2.5cm, left=2.5cm, right=2.5cm]{geometry}
\usepackage{fontspec}
\usepackage[english, bidi=basic, provide=*]{babel}

% Set default font to Noto Sans for clean academic legibility
\babelfont{rm}{Noto Sans}

% --- AMS Mathematics Packages ---
\usepackage{amsmath, amssymb, amsthm, mathtools}
\usepackage{enumitem}
\usepackage{booktabs}
\usepackage{microtype}

% --- Color Palette Definition (Green Academic Theme) ---
\usepackage[dvipsnames]{xcolor}
\definecolor{DeepEmerald}{RGB}{12, 74, 46}
\definecolor{ForestGreen}{RGB}{28, 107, 63}
\definecolor{SageGreen}{RGB}{76, 140, 102}
\definecolor{MintAccent}{RGB}{142, 202, 163}
\definecolor{LightSage}{RGB}{242, 247, 244}
\definecolor{BorderGreen}{RGB}{180, 210, 192}
\definecolor{DarkText}{RGB}{25, 38, 30}
\definecolor{LinkGreen}{RGB}{18, 90, 52}

% --- Graphics & Diagram Packages ---
\usepackage{tikz}
\usepackage{pgfplots}
\pgfplotsset{compat=1.18}
\usetikzlibrary{shapes.geometric, arrows.meta, decorations.pathmorphing, backgrounds, positioning, fit, calc, patterns}

% --- Callout Boxes Setup (tcolorbox) ---
\usepackage[most]{tcolorbox}

% Definition Box
\newtcolorbox{definitionbox}[1][]{
    enhanced,
    colback=LightSage,
    colframe=DeepEmerald,
    coltitle=white,
    fonttitle=\bfseries\sffamily,
    title={Definition: #1},
    arc=3mm,
    boxrule=0.8pt,
    drop shadow=black!10!white,
    left=4mm, right=4mm, top=3mm, bottom=3mm
}

% Theorem Box
\newtcolorbox{theorembox}[1][]{
    enhanced,
    colback=white,
    colframe=ForestGreen,
    coltitle=white,
    fonttitle=\bfseries\sffamily,
    title={Theorem: #1},
    arc=2mm,
    boxrule=1.2pt,
    borderline={0.5pt}{2pt}{ForestGreen!40},
    left=4mm, right=4mm, top=3mm, bottom=3mm
}

% Example Box
\newtcolorbox{examplebox}[1][]{
    enhanced,
    colback=LightSage!60!white,
    colframe=SageGreen,
    coltitle=white,
    fonttitle=\bfseries\sffamily,
    title={Worked Example: #1},
    arc=3mm,
    boxrule=0.8pt,
    left=4mm, right=4mm, top=3mm, bottom=3mm
}

% Category / Summary Box
\newtcolorbox{summarybox}[1][]{
    enhanced,
    colback=LightSage,
    colframe=BorderGreen,
    coltitle=DeepEmerald,
    fonttitle=\bfseries\sffamily,
    title={#1},
    arc=4mm,
    boxrule=1pt,
    left=4mm, right=4mm, top=3mm, bottom=3mm
}

% --- Custom Chapter & Section Styling ---
\usepackage{titlesec}
\titleformat{\chapter}[display]
  {\normalfont\bfseries\sffamily\color{DeepEmerald}}
  {\Large\color{ForestGreen}\chaptertitlename\ \thechapter}
  {10pt}{\Huge}
\titleformat{\section}
  {\normalfont\Large\bfseries\sffamily\color{DeepEmerald}}
  {\thesection}{1em}{}
\titleformat{\subsection}
  {\normalfont\large\bfseries\sffamily\color{ForestGreen}}
  {\thesubsection}{1em}{}

% --- Hyperlinks Configuration ---
\usepackage{hyperref}
\hypersetup{
    colorlinks=true,
    linkcolor=LinkGreen,
    citecolor=LinkGreen,
    urlcolor=LinkGreen,
    pdfauthor={Advanced Calculus Editorial Board},
    pdftitle={Advanced Calculus: Differential Equations}
}

\setlength{\parindent}{0pt}
\setlength{\parskip}{0.5em}

\begin{document}

% =========================================================
% BOOK COVER PAGE
% =========================================================
\begin{titlepage}
\begin{tikzpicture}[remember picture, overlay]
    % Background Gradient Fill
    \fill[DeepEmerald] (current page.north west) rectangle (current page.south east);
    
    % Secondary Aesthetic Panel
    \fill[ForestGreen] ($(current page.north west)+(0,-4)$) -- ($(current page.north east)+(0,-2)$) -- ($(current page.south east)+(0,5)$) -- ($(current page.south west)+(0,3)$) -- cycle;

    % Geometric Parametric Overlay Lines
    \draw[MintAccent!30, thick, domain=-10:10, samples=100, smooth] 
        plot ({2*sin(\x r) + \x/2}, {3*cos(2*\x r) - \x});
    \draw[MintAccent!20, thick, domain=-8:8, samples=100, smooth] 
        plot ({\x}, {0.2*\x*\x - 5});
    \draw[MintAccent!15, thick, domain=-10:10, samples=100, smooth] 
        plot ({\x}, {-0.1*\x*\x + 6});
        
    % Decorative Frames
    \node[anchor=north west] at ($(current page.north west)+(1.8,-2)$) {
        \textcolor{MintAccent}{\rule{4cm}{2pt}}
    };
    
    % Book Header Title Box
    \node[anchor=center] at ($(current page.center)+(0,5)$) {
        \begin{tcolorbox}[
            width=0.88\textwidth,
            colback=white,
            colframe=MintAccent,
            arc=4mm,
            boxrule=2pt,
            halign=center,
            shadow={2mm}{-2mm}{0mm}{black!30}
        ]
            \vspace{0.4cm}
            {\Huge\bfseries\sffamily\textcolor{DeepEmerald}{ADVANCED CALCULUS}}\\[0.3cm]
            {\Large\sffamily\textcolor{ForestGreen}{A Comprehensive Treatise on Differential Equations}}\\[0.2cm]
            {\small\sffamily\textcolor{DarkText}{Ordinary, Partial, and Higher-Order Linear Systems}}
            \vspace{0.4cm}
        \end{tcolorbox}
    };

    % Decorative Math Plot in Center Cover
    \node[anchor=center] at ($(current page.center)+(0,-1.5)$) {
        \begin{tikzpicture}[scale=0.85]
            % Orthogonal Field Rays on Cover
            \foreach \r in {0.8, 1.6, 2.4, 3.2} {
                \draw[MintAccent!60, thick] (0,0) circle (\r);
            }
            \foreach \angle in {0, 30, 60, 90, 120, 150, 180, 210, 240, 270, 300, 330} {
                \draw[white!70, thick, ->] (0,0) -- (\angle:3.6);
            }
            \node[fill=MintAccent, circle, inner sep=2pt] at (0,0) {};
        \end{tikzpicture}
    };

    % Author and Institution Footer Box
    \node[anchor=south] at ($(current page.south)+(0,2.5)$) {
        \begin{tcolorbox}[
            width=0.75\textwidth,
            colback=DeepEmerald!90!black,
            colframe=MintAccent!50,
            arc=3mm,
            halign=center
        ]
            {\textcolor{white}{\bfseries\sffamily Academic Lecture Notes \& Theoretical Reference}}\\[0.1cm]
            {\small\textcolor{MintAccent}{\sffamily Typeset with Publication-Quality LaTeX, TikZ, and AMS Environments}}
        \end{tcolorbox}
    };
\end{tikzpicture}
\end{titlepage}

\newpage
\tableofcontents
\newpage

% =========================================================
% CHAPTER 1: FUNDAMENTALS OF DIFFERENTIAL EQUATIONS
% =========================================================
\chapter{Fundamentals of Differential Equations}

\section{Definitions and Categorization}

Differential equations form the foundational language of dynamic modeling across pure mathematics, physics, engineering, and mathematical biology.

\begin{definitionbox}[Differential Equation (DE)]
A \textbf{differential equation} is any equation involving one or more dependent variables, one or more independent variables, and the ordinary or partial derivatives of the dependent variables with respect to the independent variables.
\end{definitionbox}

Differential equations are broadly classified into two main categories based on the structure of their derivatives:

\begin{definitionbox}[Ordinary Differential Equation (ODE)]
An \textbf{ordinary differential equation (ODE)} is a differential equation containing ordinary derivatives of one or more dependent variables with respect to a \emph{single} independent variable.
\end{definitionbox}

\begin{definitionbox}[Partial Differential Equation (PDE)]
A \textbf{partial differential equation (PDE)} is a differential equation containing partial derivatives of one or more dependent variables with respect to \emph{two or more} independent variables.
\end{definitionbox}

\begin{summarybox}[Classification of Differential Equations by Category]
    \textbf{Ordinary Differential Equations (ODE):}
    \begin{equation}
        \frac{d^2 y}{dx^2} + x y \left(\frac{dy}{dx}\right)^2 = 0
    \end{equation}
    \begin{equation}
        \frac{d^4 x}{dt^4} + 5\frac{d^2 x}{dt^2} + 3x = \sin(t)
    \end{equation}

    \textbf{Partial Differential Equations (PDE):}
    \begin{equation}
        \frac{\partial v}{\partial \sigma} + \frac{\partial v}{\partial t} = v
    \end{equation}
    \begin{equation}
        \frac{\partial^2 u}{\partial x^2} + \frac{\partial^2 u}{\partial y^2} + \frac{\partial^2 u}{\partial z^2} = 0 \quad \text{(Laplace's Equation)}
    \end{equation}

    \textbf{Standard Calculus Derivatives (Not DEs):}
    \begin{equation}
        \frac{d}{dx}\left(e^{ax}\right) = a e^{ax}, \qquad \frac{d}{dx}(u v) = u \frac{dv}{dx} + v \frac{du}{dx}
    \end{equation}
\end{summarybox}

\section{Order and Degree of a Differential Equation}

\begin{definitionbox}[Order of a Differential Equation]
The \textbf{order} of a differential equation is defined as the order of the highest-order derivative appearing in the equation.
\end{definitionbox}

\begin{definitionbox}[Degree of a Differential Equation]
The \textbf{degree} of a differential equation is the algebraic degree (highest polynomial power) of the highest-order derivative present in the equation, after the equation has been cleared of radicals and fractional exponents with respect to derivatives.
\end{definitionbox}

\begin{summarybox}[Examples of Order and Degree Classification]
    \begin{itemize}[leftmargin=0.6cm]
        \item \textbf{1st Order, 1st Degree ODE:}
        \[ \frac{dy}{dx} + y = e^x \]
        \item \textbf{2nd Order, 1st Degree ODE:}
        \[ \frac{d^2 y}{dx^2} + 3\frac{dy}{dx} + 2y = 0 \]
        \item \textbf{2nd Order, 2nd Degree ODE:}
        \[ \left(\frac{d^2 y}{dx^2}\right)^2 + y \frac{dy}{dx} = 0 \]
        \item \textbf{1st Order, 1st Degree PDE:}
        \[ \frac{\partial u}{\partial t} + \frac{\partial u}{\partial x} = 0 \]
        \item \textbf{2nd Order, 1st Degree PDE (Wave Equation):}
        \[ \frac{\partial^2 u}{\partial t^2} - c^2 \frac{\partial^2 u}{\partial x^2} = 0 \]
        \item \textbf{2nd Order, 2nd Degree PDE:}
        \[ \left(\frac{\partial^2 u}{\partial x^2}\right)^2 + \left(\frac{\partial^2 u}{\partial y^2}\right)^2 = 0 \]
    \end{itemize}
\end{summarybox}

\section{Linearity Criteria and Classification}

Linearity plays a fundamental role in solution existence, superposition properties, and analytical tractability.

\begin{definitionbox}[Linear Ordinary Differential Equation]
An $n$-th order ordinary differential equation in the dependent variable $y$ and independent variable $x$ is called \textbf{linear} if it can be written in the explicit representation:
\begin{equation}
    a_0(x) \frac{d^n y}{dx^n} + a_1(x) \frac{d^{n-1} y}{dx^{n-1}} + \dots + a_{n-1}(x) \frac{dy}{dx} + a_n(x) y = f(x)
\end{equation}
or in compact notation:
\begin{equation}
    a_0(x) y^{(n)} + a_1(x) y^{(n-1)} + \dots + a_n(x) y = f(x)
\end{equation}
where $a_0(x) \neq 0$ and $a_0(x), a_1(x), \dots, a_n(x), f(x)$ are given functions of $x$.
\end{definitionbox}

\begin{summarybox}[Key Criteria for Linearity]
An ODE is strictly \textbf{linear} if and only if it satisfies three structural rules:
\begin{enumerate}[leftmargin=0.6cm]
\item The dependent variable $y$ and all its derivatives $y', y'', \dots, y^{(n)}$ occur to the \textbf{first degree only}.
\item No product terms involving $y$ and its derivatives (e.g., $y \cdot \frac{dy}{dx}$) or products of derivatives exist.
\item No non-linear transcendental functions of $y$ or its derivatives are present (e.g., $\sin(y)$, $e^y$, $\ln(y')$).
\end{enumerate}
\end{summarybox}

\begin{definitionbox}[Nonlinear Ordinary Differential Equation]
An ordinary differential equation is classified as \textbf{nonlinear} if it violates any of the three linearity criteria outlined above. The general implicit form of an $n$-th order nonlinear ODE is:
\begin{equation}
    F\left(x, y, \frac{dy}{dx}, \frac{d^2 y}{dx^2}, \dots, \frac{d^n y}{dx^n}\right) = 0
\end{equation}
\end{definitionbox}

\begin{examplebox}[Comparing Linear and Nonlinear ODEs]
    \textbf{Linear ODE with Constant Coefficients:}
    \[ \frac{d^2 y}{dx^2} + 3\frac{dy}{dx} + 2y = 0 \]

    \textbf{Linear ODE with Variable Coefficients:}
    \[ \frac{d^4 y}{dx^4} + x^{30} \frac{dy}{dx} + 2y = x e^x \]

    \textbf{Nonlinear ODE Examples:}
    \begin{align*}
        \left(\frac{d^2 y}{dx^2}\right)^2 + y \frac{dy}{dx} &= 0 \quad \text{(Violates degree and product rule)} \\
        \frac{dy}{dx} + \sin(y) &= 0 \quad \text{(Violates transcendental function rule)} \\
        A \frac{d^2 y}{dx^2} + K \frac{d^3 y}{dx^3} + x \frac{dy}{dx} &= x e^x \frac{dx}{dy} \quad \text{(Violates derivative exponent rule)}
    \end{align*}
\end{examplebox}

\section{Origins and Engineering Applications}

Differential equations naturally model systems undergoing continuous change across science and technology:
\begin{itemize}[leftmargin=0.6cm]
    \item \textbf{Population Dynamics:} Malthusian growth models, logistic population growth, and Lotka-Volterra predator-prey systems.
    \item \textbf{Dynamical Systems \& Chaos:} Nonlinear pendulum motion, forced oscillators, and atmospheric convection.
    \item \textbf{Fluid Dynamics:} Navier-Stokes equations governing incompressible and compressible flow fields.
    \item \textbf{Heat Transfer:} Heterogeneous conduction, thermal diffusion, and transient cooling.
    \item \textbf{Mathematical Biology \& Biophysics:} Enzyme kinetics, nerve impulse propagation (Hodgkin-Huxley model), and epidemic spreads.
\end{itemize}

\section{Formation of Differential Equations}

A differential equation representing a family of curves is generated by differentiating the general equation containing arbitrary constants and subsequently eliminating those constants.

\begin{examplebox}[Elimination of Constants]
Form the ordinary differential equation corresponding to the family of curves given by:
\begin{equation}
    y = a x + b x^2 \label{eq:form1}
\end{equation}
where $a$ and $b$ are arbitrary constants.

\textbf{Solution:}
Since the equation contains two arbitrary constants ($a$ and $b$), we differentiate Equation \eqref{eq:form1} twice consecutively with respect to $x$.

\textbf{Step 1 (First Differentiation):}
\begin{equation}
    \frac{dy}{dx} = a + 2bx \label{eq:form2}
\end{equation}

\textbf{Step 2 (Second Differentiation):}
\begin{equation}
    \frac{d^2 y}{dx^2} = 2b \implies b = \frac{1}{2} \frac{d^2 y}{dx^2} \label{eq:b_expr}
\end{equation}

\textbf{Step 3 (Eliminating Constant $a$):}
Substitute the expression for $b$ from Equation \eqref{eq:b_expr} into Equation \eqref{eq:form2}:
\[
\frac{dy}{dx} = a + 2\left(\frac{1}{2} \frac{d^2 y}{dx^2}\right)x = a + x \frac{d^2 y}{dx^2}
\]
Solving for $a$:
\begin{equation}
    a = \frac{dy}{dx} - x \frac{d^2 y}{dx^2} \label{eq:a_expr}
\end{equation}

\textbf{Step 4 (Final Substitution into Original Equation):}
Substitute $a$ and $b$ back into Equation \eqref{eq:form1}:
\[
y = \left(\frac{dy}{dx} - x \frac{d^2 y}{dx^2}\right)x + \left(\frac{1}{2} \frac{d^2 y}{dx^2}\right)x^2
\]
Expanding and simplifying:
\[
y = x \frac{dy}{dx} - x^2 \frac{d^2 y}{dx^2} + \frac{1}{2}x^2 \frac{d^2 y}{dx^2} = x \frac{dy}{dx} - \frac{1}{2}x^2 \frac{d^2 y}{dx^2}
\]
Multiplying through by $2$ and rearranging into standard order yields:
\begin{equation}
    x^2 \frac{d^2 y}{dx^2} - 2x \frac{dy}{dx} + 2y = 0
\end{equation}
This is a 2nd-order, 1st-degree linear ordinary differential equation.
\end{examplebox}

% =========================================================
% CHAPTER 2: SOLUTIONS AND EXISTENCE THEORY
% =========================================================
\chapter{Solutions and Existence Theory}

\section{General, Particular, and Singular Solutions}

\begin{definitionbox}[General Solution (Complete Primitive)]
The relation involving $n$ independent arbitrary constants that satisfies an $n$-th order ordinary differential equation is termed its \textbf{general solution} (or complete primitive).
\end{definitionbox}

\begin{definitionbox}[Particular Solution]
A \textbf{particular solution} is a solution derived from the general solution by assigning specific numerical values to the arbitrary constants, typically dictated by initial or boundary conditions.
\end{definitionbox}

\begin{summarybox}[Geometric Interpretation]
Geometrically, the general solution (primitive) represents an $n$-parameter \textbf{family of curves} satisfying the differential equation, whereas a particular solution represents a \textbf{single curve} belonging to this family.
\end{summarybox}

\begin{examplebox}[Differential Equation of a Parabolic Family]
Form the differential equation of all parabolas whose axes are parallel to the $y$-axis:
\begin{equation}
    (x - h)^2 = 4a(y - k) \label{eq:parabola}
\end{equation}
where $a, h, k$ are arbitrary constants.

\textbf{Solution:}
Differentiating Equation \eqref{eq:parabola} with respect to $x$:
\[
2(x - h) = 4a \frac{dy}{dx} \implies x - h = 2a \frac{dy}{dx}
\]
Differentiating a second time with respect to $x$:
\[
1 = 2a \frac{d^2 y}{dx^2} \implies \frac{d^2 y}{dx^2} = \frac{1}{2a}
\]
Differentiating a third time eliminates the constant $a$:
\begin{equation}
    \frac{d^3 y}{dx^3} = 0
\end{equation}
This third-order ordinary differential equation represents the entire family of parabolas.
\end{examplebox}

\section{Explicit vs. Implicit Solutions}

Consider an $n$-th order real ODE expressed as $F\left(x, y, y', \dots, y^{(n)}\right) = 0$.

\begin{definitionbox}[Explicit Solution]
A real-valued function $y = f(x)$ defined on an interval $I \subseteq \mathbb{R}$ is an \textbf{explicit solution} if:
\begin{enumerate}[leftmargin=0.6cm]
\item $f(x)$ possesses continuous derivatives up to order $n$ on $I$.
\item $F\left(x, f(x), f'(x), \dots, f^{(n)}(x)\right) = 0$ holds identically for all $x \in I$.
\end{enumerate}
\end{definitionbox}

\begin{definitionbox}[Implicit Solution]
A relation $g(x,y) = 0$ is an \textbf{implicit solution} of the ODE on $I$ if the relation implicitly defines at least one function $y = f(x)$ that is an explicit solution on $I$.
\end{definitionbox}

\begin{examplebox}[Verification of an Explicit Solution]
Show that $f(x) = 2\sin(x) + 3\cos(x)$ is an explicit solution of $y'' + y = 0$ for all $x \in \mathbb{R}$.

\textbf{Solution:}
Given $f(x) = 2\sin(x) + 3\cos(x)$, we compute its derivatives:
\[
f'(x) = 2\cos(x) - 3\sin(x), \qquad f''(x) = -2\sin(x) - 3\cos(x)
\]
Substituting into the differential equation:
\[
f''(x) + f(x) = \Big(-2\sin(x) - 3\cos(x)\Big) + \Big(2\sin(x) + 3\cos(x)\Big) = 0
\]
Since this holds identically for all $x \in \mathbb{R}$, $f(x)$ is an explicit solution.
\end{examplebox}

\section{Initial Value Problems (IVP) and Boundary Value Problems (BVP)}

\begin{definitionbox}[Initial Value Problem (IVP)]
An \textbf{Initial Value Problem (IVP)} consists of a differential equation accompanied by supplementary conditions specified at a \textbf{single point} $x_0$:
\begin{equation}
    y(x_0) = y_0, \quad y'(x_0) = y_1, \quad \dots, \quad y^{(n-1)}(x_0) = y_{n-1}
\end{equation}
\end{definitionbox}

\begin{definitionbox}[Boundary Value Problem (BVP)]
A \textbf{Boundary Value Problem (BVP)} consists of a differential equation accompanied by supplementary conditions specified at \textbf{two or more distinct points} $x_a, x_b$:
\begin{equation}
    y(x_a) = y_a, \qquad y(x_b) = y_b
\end{equation}
\end{definitionbox}

\begin{examplebox}[Contrasting IVP and BVP]
    \textbf{Initial Value Problem (IVP):}
    \[ \frac{d^2 y}{dx^2} + y = 0, \qquad y(1) = 3, \quad y'(1) = -4 \quad \text{(Both conditions at $x = 1$)} \]

    \textbf{Boundary Value Problem (BVP):}
    \[ \frac{d^2 y}{dx^2} + y = 0, \qquad y(0) = 1, \quad y\left(\frac{\pi}{2}\right) = 5 \quad \text{(Conditions at $x = 0$ and $x = \pi/2$)} \]
\end{examplebox}

\section{Picard-Lindelöf Existence and Uniqueness Theorem}

\begin{theorembox}[Picard-Lindelöf Theorem]
Consider the first-order initial value problem:
\begin{equation}
    \frac{dy}{dx} = f(x,y), \qquad y(x_0) = y_0
\end{equation}
Let $f(x,y)$ and its partial derivative $\frac{\partial f}{\partial y}$ be continuous in a closed rectangular region $D = \{(x,y) \mid |x - x_0| \le a, |y - y_0| \le b\}$.

Then, there exists a \textbf{unique solution} $y = \phi(x)$ defined on an interval $|x - x_0| \le h$ (where $h \le a$), satisfying the initial condition $\phi(x_0) = y_0$.
\end{theorembox}

% =========================================================
% CHAPTER 3: FIRST-ORDER ORDINARY DIFFERENTIAL EQUATIONS
% =========================================================
\chapter{First-Order Ordinary Differential Equations}

A first-order ODE can be expressed in derivative form $\frac{dy}{dx} = f(x,y)$ or differential form $M(x,y)dx + N(x,y)dy = 0$.

\section{Separable Equations and Reducible Forms}

\begin{definitionbox}[Separable Differential Equation]
A differential equation is \textbf{separable} if it can be written as:
\begin{equation}
    F(x)G(y)dx + f(x)g(y)dy = 0 \implies \frac{F(x)}{f(x)}dx + \frac{g(y)}{G(y)}dy = 0
\end{equation}
\end{definitionbox}

\begin{examplebox}[Separable Equation Solution]
Solve the initial value problem:
\[ x \ln(y) dx + (x^2 + 1) \cos(y) dy = 0, \qquad y(1) = \frac{\pi}{2} \]

\textbf{Solution:}
Separating variables:
\[ \frac{x}{x^2 + 1} dx + \frac{\cos(y)}{\ln(y)} dy = 0 \]
For the initial condition $y(1) = \frac{\pi}{2}$, integrating both terms yields the complete explicit curve passing through $(1, \pi/2)$.
\end{examplebox}

\begin{examplebox}[Reducible to Separable Form]
Solve $\frac{dy}{dx} = (4x + y + 1)^2$.

\textbf{Solution:}
Let $u = 4x + y + 1$. Differentiating with respect to $x$:
\[ \frac{du}{dx} = 4 + \frac{dy}{dx} \implies \frac{dy}{dx} = \frac{du}{dx} - 4 \]
Substituting into the ODE:
\[ \frac{du}{dx} - 4 = u^2 \implies \frac{du}{u^2 + 4} = dx \]
Integrating both sides:
\[ \frac{1}{2} \arctan\left(\frac{u}{2}\right) = x + C \implies \arctan\left(\frac{4x + y + 1}{2}\right) = 2x + 2C \]
Taking the tangent of both sides:
\begin{equation}
    4x + y + 1 = 2 \tan(2x + K)
\end{equation}
\end{examplebox}

\section{Homogeneous Differential Equations}

\begin{definitionbox}[Homogeneous Function and DE]
A function $f(x,y)$ is homogeneous of degree $n$ if $f(k x, k y) = k^n f(x,y)$. The equation $M(x,y)dx + N(x,y)dy = 0$ is \textbf{homogeneous} if $M$ and $N$ are homogeneous functions of the same degree.
\end{definitionbox}

Substitution $y = v x$ (hence $\frac{dy}{dx} = v + x \frac{dv}{dx}$) reduces any homogeneous equation into a separable equation.

\begin{examplebox}[Homogeneous Differential Equation]
Solve $(x^2 + y^2)dx + 2xy dy = 0$.

\textbf{Solution:}
Expressing in derivative form:
\[ \frac{dy}{dx} = -\frac{x^2 + y^2}{2xy} \]
Substituting $y = v x$ and $\frac{dy}{dx} = v + x \frac{dv}{dx}$:
\[ v + x \frac{dv}{dx} = -\frac{x^2 + v^2 x^2}{2 x^2 v} = -\frac{1 + v^2}{2v} \]
Rearranging:
\[ x \frac{dv}{dx} = -\frac{1 + v^2}{2v} - v = -\frac{1 + 3v^2}{2v} \implies \frac{2v}{1 + 3v^2} dv = -\frac{dx}{x} \]
Integrating both sides:
\[ \frac{1}{3} \ln(1 + 3v^2) = -\ln(x) + \ln(C) \implies \ln\Big(x^3(1 + 3v^2)\Big) = \ln(C') \]
Substituting $v = y/x$:
\begin{equation}
    x^3 \left(1 + 3\frac{y^2}{x^2}\right) = C \implies x^3 + 3x y^2 = C
\end{equation}
\end{examplebox}

\section{First-Order Linear Differential Equations}

A first-order linear ODE has the canonical form:
\begin{equation}
    \frac{dy}{dx} + P(x)y = Q(x)
\end{equation}
Multiplying by the \textbf{Integrating Factor (I.F.)} $\mu(x) = e^{\int P(x)dx}$ transforms the left-hand side into an exact product derivative:
\begin{equation}
    \frac{d}{dx}\left[ y e^{\int P(x)dx} \right] = Q(x) e^{\int P(x)dx} \implies y = e^{-\int P(x)dx} \left[ \int Q(x) e^{\int P(x)dx} dx + C \right]
\end{equation}

\section{Bernoulli Equations}

\begin{definitionbox}[Bernoulli Equation]
An equation of the form:
\begin{equation}
    \frac{dy}{dx} + P(x)y = Q(x)y^n \qquad (n \neq 0, 1)
\end{equation}
is a \textbf{Bernoulli differential equation}.
\end{definitionbox}

Dividing by $y^n$ and substituting $v = y^{1-n}$ transforms it into a first-order linear ODE in $v$:
\begin{equation}
    \frac{dv}{dx} + (1-n)P(x)v = (1-n)Q(x)
\end{equation}

\section{Exact Differential Equations}

\begin{theorembox}[Condition for Exactness]
The differential form $M(x,y)dx + N(x,y)dy = 0$ is \textbf{exact} if and only if:
\begin{equation}
    \frac{\partial M}{\partial y} = \frac{\partial N}{\partial x}
\end{equation}
If exact, there exists a potential function $F(x,y)$ such that $\frac{\partial F}{\partial x} = M$ and $\frac{\partial F}{\partial y} = N$, giving the general solution $F(x,y) = C$.
\end{theorembox}

\section{Orthogonal Trajectories}

An \textbf{orthogonal trajectory} is a curve that intersects every member of a given family of curves at right angles ($90^\circ$).

\begin{summarybox}[Rule for Finding Orthogonal Trajectories]
\begin{enumerate}[leftmargin=0.6cm]
\item For Cartesian family $f(x,y,c) = 0$: find ODE $F(x,y,y') = 0$, replace $y' \to -\frac{1}{y'}$, and integrate.
\item For Polar family $f(r,\theta,c) = 0$: find ODE $F(r,\theta,r') = 0$, replace $\frac{dr}{d\theta} \to -r^2 \frac{d\theta}{dr}$, and integrate.
\end{enumerate}
\end{summarybox}

\begin{figure}[htbp]
\centering
\begin{tikzpicture}[scale=1.1]
    % Concentric Circles Family (Original)
    \foreach \r in {0.6, 1.2, 1.8, 2.4, 3.0} {
        \draw[ForestGreen, thick] (0,0) circle (\r);
    }
    % Orthogonal Lines Family
    \foreach \angle in {0, 30, 60, 90, 120, 150, 180, 210, 240, 270, 300, 330} {
        \draw[DeepEmerald!80, dashed, thick] (-\angle:3.3) -- (\angle:3.3);
    }
    % Axes
    \draw[->, >=stealth, thick] (-3.6,0) -- (3.6,0) node[right] {$x$};
    \draw[->, >=stealth, thick] (0,-3.6) -- (0,3.6) node[above] {$y$};
    \node[fill=white, inner sep=1pt] at (2.2,2.2) {\textcolor{ForestGreen}{$x^2 + y^2 = C^2$}};
    \node[fill=white, inner sep=1pt] at (-2.5,1.2) {\textcolor{DeepEmerald}{$y = k x$}};
\end{tikzpicture}
\caption{Orthogonal trajectories: Concentric circles $x^2 + y^2 = C^2$ (solid green) intersected perpendicularly by radial lines $y = k x$ (dashed emerald).}
\label{fig:orthogonal_circles}
\end{figure}

\begin{examplebox}[Parabolic Orthogonal Trajectories]
Find the orthogonal trajectories of the parabolic family $y = c x^2$.

\textbf{Solution:}
Differentiating $y = c x^2$ gives $\frac{dy}{dx} = 2c x = 2\left(\frac{y}{x^2}\right)x = \frac{2y}{x}$.
Replacing $\frac{dy}{dx}$ with $-\frac{dx}{dy}$ for orthogonality:
\[ -\frac{dx}{dy} = \frac{2y}{x} \implies x dx + 2y dy = 0 \]
Integrating both sides:
\begin{equation}
    \frac{x^2}{2} + y^2 = C' \implies x^2 + 2y^2 = K^2
\end{equation}
The orthogonal trajectories form a family of \textbf{ellipses} centered at the origin.
\end{examplebox}

\begin{figure}[htbp]
\centering
\begin{tikzpicture}[scale=0.95]
    \begin{axis}[
        xmin=-3, xmax=3, ymin=-3, ymax=3,
        axis lines=middle,
        xlabel={$x$}, ylabel={$y$},
        grid=both,
        grid style={line width=.1pt, draw=gray!20},
        major grid style={line width=.2pt, draw=gray!40},
        domain=-2.5:2.5,
        samples=100
    ]
        % Parabolas
        \addplot[ForestGreen, thick] {0.4*x^2};
        \addplot[ForestGreen, thick] {0.8*x^2};
        \addplot[ForestGreen, thick] {-0.4*x^2};
        \addplot[ForestGreen, thick] {-0.8*x^2};

        % Ellipses (Orthogonal trajectories)
        \addplot[DeepEmerald, dashed, thick, domain=-1.414:1.414] {sqrt(1 - 0.5*x^2)};
        \addplot[DeepEmerald, dashed, thick, domain=-1.414:1.414] {-sqrt(1 - 0.5*x^2)};
        \addplot[DeepEmerald, dashed, thick, domain=-2.0:2.0] {sqrt(2 - 0.5*x^2)};
        \addplot[DeepEmerald, dashed, thick, domain=-2.0:2.0] {-sqrt(2 - 0.5*x^2)};
    \end{axis}
\end{tikzpicture}
\caption{Parabolas $y = c x^2$ (solid green) and their orthogonal elliptical trajectories $x^2 + 2y^2 = K^2$ (dashed dark green).}
\label{fig:parabola_ellipse}
\end{figure}

% =========================================================
% CHAPTER 4: HIGHER-ORDER LINEAR DIFFERENTIAL EQUATIONS
% =========================================================
\chapter{Higher-Order Linear Differential Equations}

\section{General Theory and Superposition Principle}

An $n$-th order non-homogeneous linear ODE is represented by:
\begin{equation}
    a_0(x)\frac{d^n y}{dx^n} + a_1(x)\frac{d^{n-1} y}{dx^{n-1}} + \dots + a_n(x)y = F(x)
\end{equation}

\begin{theorembox}[Superposition Principle]
If $y_1, y_2, \dots, y_m$ are solutions of the homogeneous linear ODE ($F(x) = 0$), then any linear combination:
\begin{equation}
    y(x) = c_1 y_1(x) + c_2 y_2(x) + \dots + c_m y_m(x)
\end{equation}
is also a solution.
\end{theorembox}

\section{Linear Independence and the Wronskian Determinant}

\begin{definitionbox}[Wronskian Determinant]
For $n$ differentiable functions $f_1, f_2, \dots, f_n$, the \textbf{Wronskian} $W(f_1, \dots, f_n)$ is defined as:
\begin{equation}
    W(f_1, f_2, \dots, f_n)(x) = 
    \begin{vmatrix}
        f_1(x) & f_2(x) & \dots & f_n(x) \\
        f_1'(x) & f_2'(x) & \dots & f_n'(x) \\
        \vdots & \vdots & \ddots & \vdots \\
        f_1^{(n-1)}(x) & f_2^{(n-1)}(x) & \dots & f_n^{(n-1)}(x)
    \end{vmatrix}
\end{equation}
\end{definitionbox}

\begin{theorembox}[Wronskian Test for Linear Independence]
Solutions $y_1, y_2, \dots, y_n$ of an $n$-th order homogeneous linear ODE are \textbf{linearly independent} on an interval $I$ if and only if $W(y_1, y_2, \dots, y_n)(x) \neq 0$ for at least one $x \in I$.
\end{theorembox}

\begin{examplebox}[Wronskian Verification]
Verify that $e^x, e^{-x}, e^{2x}$ are linearly independent solutions of $y''' - 2y'' - y' + 2y = 0$.

\textbf{Solution:}
Evaluating the Wronskian:
\[
W(e^x, e^{-x}, e^{2x}) = 
\begin{vmatrix}
    e^x & e^{-x} & e^{2x} \\
    e^x & -e^{-x} & 2e^{2x} \\
    e^x & e^{-x} & 4e^{2x}
\end{vmatrix}
= e^x \cdot e^{-x} \cdot e^{2x}
\begin{vmatrix}
    1 & 1 & 1 \\
    1 & -1 & 2 \\
    1 & 1 & 4
\end{vmatrix}
\]
Computing the determinant:
\[
1(-4 - 2) - 1(4 - 2) + 1(1 - (-1)) = -6 - 2 + 2 = -6
\]
Thus $W = -6 e^{2x} \neq 0$ for all $x \in \mathbb{R}$, proving linear independence.
\end{examplebox}

\section{Reduction of Order}

If one non-trivial solution $y_1(x)$ of an $n$-th order linear homogeneous ODE is known, the substitution $y(x) = v(x) y_1(x)$ reduces the ODE to an $(n-1)$-th order linear ODE in $w = v'$.

\section{Homogeneous Linear ODEs with Constant Coefficients}

For $a_0 y^{(n)} + a_1 y^{(n-1)} + \dots + a_n y = 0$, substituting $y = e^{m x}$ yields the \textbf{Auxiliary (Characteristic) Equation}:
\begin{equation}
    a_0 m^n + a_1 m^{n-1} + \dots + a_n = 0
\end{equation}

\begin{summarybox}[Root Cases for the Auxiliary Equation]
\begin{enumerate}[leftmargin=0.6cm]
\item \textbf{Distinct Real Roots} ($m_1 \neq m_2 \dots \neq m_n$):
    \[ y_c = c_1 e^{m_1 x} + c_2 e^{m_2 x} + \dots + c_n e^{m_n x} \]
\item \textbf{Repeated Real Roots} ($m_1 = m_2 = \dots = m_k = m$):
    \[ y_c = (c_1 + c_2 x + c_3 x^2 + \dots + c_k x^{k-1}) e^{m x} \]
\item \textbf{Complex Conjugate Roots} ($m = \alpha \pm i \beta$):
    \[ y_c = e^{\alpha x} \Big( c_1 \cos(\beta x) + c_2 \sin(\beta x) \Big) \]
\end{enumerate}
\end{summarybox}

\begin{examplebox}[Complex Roots Initial Value Problem]
Solve the IVP: $y'' - 6y' + 25y = 0$, with $y(0) = -3, y'(0) = -1$.

\textbf{Solution:}
Auxiliary equation: $m^2 - 6m + 25 = 0 \implies m = \frac{6 \pm \sqrt{36 - 100}}{2} = 3 \pm 4i$.
General solution:
\[ y(x) = e^{3x} \Big( c_1 \cos(4x) + c_2 \sin(4x) \Big) \]
Applying $y(0) = -3 \implies c_1 = -3$.
Differentiating and applying $y'(0) = -1$:
\[ y'(x) = 3e^{3x}(c_1 \cos 4x + c_2 \sin 4x) + e^{3x}(-4c_1 \sin 4x + 4c_2 \cos 4x) \]
\[ -1 = 3(-3) + 4c_2 \implies 4c_2 = 8 \implies c_2 = 2 \]
The unique solution is:
\begin{equation}
    y(x) = e^{3x} \Big( 2\sin(4x) - 3\cos(4x) \Big)
\end{equation}
\end{examplebox}

\section{Non-Homogeneous ODEs and Particular Integrals}

The general solution of a non-homogeneous ODE is $y = y_c + y_p$ (Complementary Function + Particular Integral).

Using the differential operator $D \equiv \frac{d}{dx}$, $y_p = \frac{1}{f(D)} F(x)$.

\begin{summarybox}[Inverse Operator $D$-Rules]
\begin{enumerate}[leftmargin=0.6cm]
\item \textbf{Exponential Source} $F(x) = e^{ax}$:
    \[ \frac{1}{f(D)} e^{ax} = \frac{1}{f(a)} e^{ax} \quad \text{if } f(a) \neq 0 \]
    \textit{Failure Case:} If $f(a) = 0$, $\frac{1}{f(D)} e^{ax} = x \frac{1}{f'(a)} e^{ax}$.
\item \textbf{Sinusoidal Source} $F(x) = \sin(ax)$ or $\cos(ax)$: Replace $D^2 \to -a^2$.
\item \textbf{Polynomial Source} $F(x) = x^m$: Expand $[f(D)]^{-1}$ as a binomial series in ascending powers of $D$.
\item \textbf{Exponential Shift}: $\frac{1}{f(D)}\left( e^{ax} V(x) \right) = e^{ax} \frac{1}{f(D+a)} V(x)$.
\end{enumerate}
\end{summarybox}

\end{document}
